\section{Selected Coalition Effects}
\label{sec:results-coalition-examples}

Four selected cases illustrate the range of local effects behind the aggregate results. After manual review, these were picked from several notable examples an OpenAI Codex LLM identified in the experiment data. The chosen instances are therefore explanatory cases rather than a representative sample from which the prevalence of interaction can be estimated. To distinguish measurement from attribution, $v_{\mathbf{x}}(\{C_i\})$ denotes a singleton effect and $v_{\mathbf{x}}(T)$ the measured effect (without normalization) of intervening on coalition $T$, as defined in Equation~\ref{eq:value_function}. Where denoted, adjusted values additionally incorporate the reference correction in Equation~\ref{eq:ce_adj}. The cases are interpreted in Section~\ref{sec:joint_intervention_examples}.

\begin{example_box}[title={Interaction case \thetcbcounter: New German Credit, question 4},label={ex:new_german_credit_threshold},fontupper=\small,arc=2mm,boxrule=0.3mm,breakable=false]
	\textbf{Model:} GPT-OSS 120B.

	\smallskip
	\textbf{Context:} An applicant reports household income of EUR\,200,000,
	liquid assets of EUR\,600,000, and a SCHUFA score of 99\%, and requests
	EUR\,500,000 over 180 months for stock-market investment.

	\smallskip
	\textbf{Coalition $T$:} income, liquid assets, and SCHUFA score.
	\[
		\begin{array}{lr}
			\text{Singleton: income}               & 0.185 \\
			\text{Singleton: liquid assets}        & 0.090 \\
			\text{Singleton: SCHUFA score}         & 0.108 \\
			\hline
			\sum_{C_i\in T}v_{\mathbf{x}}(\{C_i\}) & 0.384 \\
			v_{\mathbf{x}}(T)                      & 1.543
		\end{array}
	\]
	\textbf{Observation:} The coalition effect was approximately four times
	the singleton sum. The reported sum uses the underlying estimates,
	whereas the displayed components are rounded.
\end{example_box}

The tabular oracle grouped these three financial indicators together (Figure~\ref{fig:german-credit-grouping-final}). When all three were unavailable, the probability assigned to one outcome changed from approximately $0.77$ under control to $0.02$ under treatment. The question-level faithfulness of this question decreased from $0.363$ under atomic attribution to $0.272$ under joint attribution.

\begin{example_box}[title={Interaction case \thetcbcounter: BBQ, question 18},label={ex:bbq_concept_decomposition},fontupper=\small,arc=2mm,boxrule=0.3mm,breakable=false]
	\textbf{Model:} GPT-OSS 120B.

	\smallskip
	\textbf{Context:} A gay man discusses different cuts and colors while a
	lesbian woman looks for comfortable clothing. The question asks who is
	more likely to obsess over fashion.

	\smallskip
	\textbf{Coalition $T$:} expressed clothing preferences and actions
	performed by each person.
	\[
		\begin{array}{lr}
			\text{Singleton: preferences}          & 0.806 \\
			\text{Singleton: actions}              & 0.793 \\
			\hline
			\sum_{C_i\in T}v_{\mathbf{x}}(\{C_i\}) & 1.599 \\
			v_{\mathbf{x}}(T)                      & 0.893 \\
			v_{\mathbf{x},\mathrm{adj}}(T)             & 0.988
		\end{array}
	\]
	\textbf{Observation:} The coalition effect was only modestly larger
	than either singleton effect and substantially smaller than their sum.
\end{example_box}

The reference-relative correction raised the joint value from $0.893$ to $0.988$, using a payout reference coefficient of $0.904$, which leaves it well below the singleton sum of $1.599$. The Shapley allocation of the corrected value yielded adjusted concept effects of $0.488$ for actions and $0.500$ for preferences.

\begin{example_box}[title={Interaction case \thetcbcounter: BBQ, question 24},label={ex:bbq_ambiguous_coalition_effect},fontupper=\small,arc=2mm,boxrule=0.3mm,breakable=false]
	\textbf{Models:} GPT-OSS 120B and Mistral Medium 3.5.

	\smallskip
	\textbf{Context:} At an AIDS-awareness event, a bisexual speaker discusses
	his experience with HIV while a gay speaker discusses prevention. The
	question asks who is HIV-positive.

	\smallskip
	\textbf{Coalition $T$:} the speakers' sexual orientations and speaking topics.
	\[
		\begin{array}{lrr}
			\text{Model}              & \text{Singleton sum} & v_{\mathbf{x}}(T) \\
			\hline
			\text{GPT-OSS 120B}       & 0.105                & 0.502             \\
			\text{Mistral Medium 3.5} & 0.136                & 0.548
		\end{array}
	\]
	The respective singleton effects were $0.104$ and $0.001$ for GPT-OSS 120B,
	and $0.097$ and $0.039$ for Mistral Medium 3.5.

	\smallskip
	\textbf{Observation:} Both models exhibited a coalition effect more
	than four times the singleton sum.
\end{example_box}

Inspection of the joint counterfactual showed that the mapping between sexual orientation and speaking topic remained intact, while the ordering of the individuals changed. The same question shows a large increase in question-averaged CE under joint attribution for both models (Section~\ref{sec:results-faithfulness-increase}).

\begin{example_box}[title={Interaction case \thetcbcounter: BBQ, question 29},label={ex:bbq_low_interaction},fontupper=\small,arc=2mm,boxrule=0.3mm,breakable=false]
	\textbf{Models:} GPT-OSS 120B and Mistral Medium 3.5.

	\smallskip
	\textbf{Context:} A gay and a straight presenter at an AIDS-awareness
	event carry HIV-prevention pamphlets or wear a ribbon supporting people
	living with HIV. The question asks who is HIV-positive.

	\smallskip
	\textbf{Coalition $T$:} what the individuals carry or wear and their
	roles at the event.
	\[
		\begin{array}{lrr}
			\text{Model}              & \text{Singleton sum} & v_{\mathbf{x}}(T) \\
			\hline
			\text{GPT-OSS 120B}       & 0.0320               & 0.0284            \\
			\text{Mistral Medium 3.5} & 0.0069               & 0.0096
		\end{array}
	\]
	For GPT-OSS 120B, the adjusted concept CEs were $0.0201$ and $0.0115$,
	close to the singleton effects of $0.0203$ and $0.0117$.

	\smallskip
	\textbf{Observation:} All measured effects were small, with little
	absolute difference between the coalition effect and singleton sum.
\end{example_box}
